\section{A uniform finite construction for every factor count}
\label{app:uniformallk}

For each number $k$ of prime factors we give a finite construction
whose output is exactly the set of composite solutions with $k$ prime
factors, with one proof for all $k$. Finite procedures for a fixed
number of prime factors are known (see for instance~\cite{MS26}); the
point here is a version with explicit bounds and a short complete
proof. Whether the output is empty for every $k$ is Lehmer's
conjecture.

Fix $\varepsilon\in\{-1,1\}$ and consider composite solutions of
\begin{equation}\label{eq:uniformtarget}
 n+\varepsilon=M\varphi(n),\qquad M\in\mathbb Z_{\ge2}.
\end{equation}
Such a solution is odd and squarefree: evenness contradicts the parity of
the equation, and a squared prime divisor of $n$ would divide both
$\varphi(n)$ and $n+\varepsilon$. Also $\gcd(n,M\varphi(n))=1$.
Write the prime factors in increasing order. For a proper prefix put
\[
 A=\prod_{i=1}^{j}p_i,\qquad B=\prod_{i=1}^{j}(p_i-1),\qquad
 t=k-j,\qquad D=MB-A.
\]
The empty prefix has $A=B=1$ and $t=k$. If $c$ is the product of the
remaining primes, then
\begin{equation}\label{eq:uniformdefectpositive}
 D\varphi(c)=A(c-\varphi(c))+\varepsilon>0.
\end{equation}
For a nonempty prefix, $A\ge3$ and $c-\varphi(c)\ge1$; for the empty
prefix, $D=M-1\ge1$. Thus every proper prefix of a solution has $D>0$.

\begin{lemma}[Uniform bound for the next prime]\label{lem:uniformnext}
Let $p$ be the next prime of a solution extending the prefix, with $t$
primes remaining. Then
\begin{align}
 p&\le 1+\left\lfloor\frac{tA-1}{D}\right\rfloor
       &&(\varepsilon=-1),\label{eq:uniformlinearminus}\\
 p&\le 1+\left\lfloor\frac{tA+1}{D}\right\rfloor
       &&(\varepsilon=1).\label{eq:uniformlinearplus}
\end{align}
For $\varepsilon=-1$ one also has
\begin{equation}\label{eq:uniformoddcutoff}
 p\le2+\left\lfloor
 \frac{2tA^2-1}{D(MB+A)}\right\rfloor.
\end{equation}
\end{lemma}

\begin{proof}
Let $r_1=p<r_2<\cdots<r_t$ be the remaining primes. The function
$x/(x-1)$ decreases for $x>1$. Since $r_i\ge p+i-1$,
\[
 \prod_{i=1}^{t}\frac{r_i}{r_i-1}
 \le\prod_{i=1}^{t}\frac{p+i-1}{p+i-2}
 =1+\frac{t}{p-1}.
\]
For the minus sign, $MB<A\prod_i r_i/(r_i-1)$, so
$D(p-1)<tA$. All quantities are integers, giving
\eqref{eq:uniformlinearminus}. For the plus sign,
\[
 MB=A\prod_i\frac{r_i}{r_i-1}
       +\frac{1}{\prod_i(r_i-1)}
 \le A+\frac{tA}{p-1}+\frac{1}{(p-1)^t}.
\]
It follows that $D(p-1)\le tA+(p-1)^{1-t}\le tA+1$, which gives
\eqref{eq:uniformlinearplus}.

For the sharper minus-sign estimate, oddness gives
$r_i\ge p+2(i-1)$. For $x>2$,
\[
 \left(\frac{x}{x-1}\right)^2<\frac{x}{x-2},
\]
because $x(x-2)<(x-1)^2$. Multiplying the resulting inequalities yields
\[
 \left(\prod_i\frac{r_i}{r_i-1}\right)^2
 <\prod_{i=1}^{t}\frac{p+2i-2}{p+2i-4}
 =\frac{p+2t-2}{p-2}.
\]
Consequently $(MB)^2(p-2)<A^2(p+2t-2)$, or
$D(MB+A)(p-2)<2tA^2$. The integer form of this strict inequality is
\eqref{eq:uniformoddcutoff}.
\end{proof}

\begin{corollary}[Finite quotient range]\label{cor:uniformquotient}
For $k\ge2$, the quotient of a composite solution satisfies
\[
 2\le M\le
 \begin{cases}
 \min\{\lfloor\sqrt{2k}\rfloor,\lfloor(k+1)/2\rfloor\},
       &\varepsilon=-1,\\
 \lfloor(k+3)/2\rfloor,&\varepsilon=1.
 \end{cases}
\]
\end{corollary}

\begin{proof}
At the empty prefix $D=M-1$ and the first prime is at least $3$.
The linear inequalities give the displayed linear bounds. For the minus
sign, the squared estimate in the proof of
Lemma~\ref{lem:uniformnext} gives $M^2<2k+1$. Since $M^2$ is an
integer, $M^2\le2k$.
\end{proof}

\begin{theorem}[Complete finite tree for each $k$]\label{thm:uniformtree}
For every $k\ge2$ and $\varepsilon\in\{-1,1\}$, the following finite
construction outputs exactly the composite integers with $k$ distinct
prime factors satisfying \eqref{eq:uniformtarget}, each once.
\begin{enumerate}[label=\textup{(\roman*)}]
\item For each integer $M$ in Corollary~\ref{cor:uniformquotient}, start
with $(A,B,t)=(1,1,k)$ and no chosen primes.
\item At a node with $t\ge2$, let $r$ be the last chosen prime, or $1$
at the empty prefix. Enumerate the primes $p$ such that
\[
 p>r,\qquad p>\frac{MB}{D},
\]
and $p$ is at most the bound of Lemma~\ref{lem:uniformnext} (the minimum
of its two minus-sign bounds when applicable). Keep only choices with
$\gcd(A,p-1)=1$ and $p\nmid M$. Replace the state by
\[
 (A,B,t)\longmapsto(Ap,B(p-1),t-1).
\]
\item At $t=1$, calculate the single rational number
\begin{equation}\label{eq:uniformtail}
 q=\frac{MB+\varepsilon}{D}
   =1+\frac{A+\varepsilon}{D}.
\end{equation}
Output $Aq$ exactly when $q$ is an integer prime greater than $r$.
\end{enumerate}
One may also discard any node containing $3$ unless
$M\equiv-\varepsilon\pmod3$.
\end{theorem}

\begin{proof}
Every node has $D>0$ and therefore a finite upper bound for its next
prime. The condition $p>MB/D$ makes the next proper-prefix defect
$pD-MB$ positive. The depth decreases by one, so induction on $t$
proves termination. Primality itself can be decided by trial division
through the integer square root; no probable-prime assumption is needed.

Every genuine solution has a quotient in the stated finite range and a
unique increasing list of primes. Its proper prefixes satisfy
\eqref{eq:uniformdefectpositive} and Lemma~\ref{lem:uniformnext}.
The gcd tests follow from $\gcd(n,M\varphi(n))=1$, so they cannot remove
that list. Its last prime is forced by \eqref{eq:uniformtail}. This
proves completeness.

Conversely, an accepted tail is a prime different from all prefix
primes. Thus $\varphi(Aq)=B(q-1)$, and the identity
$Dq=MB+\varepsilon$, with $D=MB-A$, gives
$Aq+\varepsilon=MB(q-1)$. Every output therefore satisfies the target
equation. Its increasing factorization and its quotient are unique.

Finally, if $3\mid n$, no other prime factor can be $1$ modulo $3$,
by the same gcd condition. Hence $\varphi(n)\equiv2\pmod3$ and
$M\varphi(n)\equiv\varepsilon\pmod3$, giving
$M\equiv-\varepsilon\pmod3$. This proves the optional prune.
\end{proof}

\subsection*{Implementation}
The program \texttt{uniform\_tree.py} (in the folder \texttt{anc} of
the source) carries out the construction in integer arithmetic with
trial division. A run that reaches one of its resource limits reports
that it is incomplete and is never counted as an empty search. For the
minus sign it finishes $2\le k\le13$ with empty output, and for the
plus sign $2\le k\le6$ with the six composite solutions of
\eqref{eq:second}. It was compared with direct evaluation on all $627$
sets of two to five odd primes up to $31$, for both signs. For $k=14$
a run stops at its limit on the length of an interval.

Write $\mathcal E_{-}(k)$ for the output for the minus sign. Lehmer's
conjecture is the assertion that $\mathcal E_{-}(k)=\varnothing$ for
every $k\ge2$; the construction gives no induction on $k$.
